\documentclass[10pt,a4paper]{amsart}
\usepackage[margin=19mm]{geometry}
\usepackage{amsmath,amssymb,mathtools}
\usepackage[scr=rsfs,cal=euler]{mathalfa}
\usepackage{xfrac}
\usepackage[unicode,hidelinks,bookmarks=false]{hyperref}
\newcommand{\ie}{i.e.\ }
\newcommand{\cf}{cf.\ }
\newcommand{\resp}{respectively\ }
\newcommand{\Subsection}{\subsection}
\providecommand{\nobreakdash}{\nobreak\hbox{-}}
\setlength{\emergencystretch}{2em}
\multlinegap0pt
\usepackage{amssymb}
\usepackage{enumerate}
\usepackage{xcolor}
\usepackage{tikz}
\usepackage[shortlabels]{enumitem}
\newtheorem{theorem}{Theorem}
\newcommand\Asc{\operatorname{Asc}}
\newcommand\Des{\operatorname{Des}}
\newcommand\maj{\operatorname{maj}}
\newcommand\op{^{\rm op}}
\title{$q$-chromatic polynomials}
\author{Esme Bajo, Matthias Beck, Andr\'es R. Vindas-Mel\'endez}
\date{Source: arXiv:2403.19573v2 (CC BY 4.0)}
\begin{document}
\maketitle

\noindent\emph{Source and licence.} $q$-chromatic polynomials. Version arXiv:2403.19573v2 (CC BY 4.0), distributed by arXiv under the Creative Commons Attribution 4.0 International licence, http://creativecommons.org/licenses/by/4.0/, as recorded in the arXiv licence metadata for this exact version and shown on the abstract page https://arxiv.org/abs/2403.19573v2. Full licence notice: \texttt{licenses/arxiv-2403.19573-CC-BY-4.0.txt}.\par\smallskip

\noindent\textit{Editorial scope.} Counted unit: the article's theorem theorem (source0.tex lines 1008--1016). The article's complete original proof of it is delivered with it (source0.tex lines 1018--1046, one \texttt{proof} environment of 29 lines). No mathematics is written, repaired or re-derived: the statement and the proof are the article's own lines, in the article's own notation, and no retained line is substituted or re-flowed.  Retained uncounted as the source-local context the proof needs: lines 980--984.  Nothing was omitted from any retained span, so the package carries no omission record.  No retained line contains a citation, so the wrapper carries no bibliography and no bibliography entry.  No retained line contains a figure environment, an image-inclusion command or drawing code, so no image is delivered and the package declares no asset dependency.  Editorial formatting only, in addition: an AMS \texttt{amsart} preamble that restates the article's own declarations and macros used by the retained lines, namely the environments \texttt{theorem} on separate counters, together with the standard packages those lines need.  The retained environments print in this document as Theorem 1 in order of appearance, whereas the article prints the same items with the numbering of its own full document.  The complete native source of the article as distributed in the arXiv e-print for this exact version is bundled unchanged as \texttt{sources/156849/source0.tex}.
\begin{equation}\label{eq:ppartgenfct}
  P_\Pi(q)
  \ = \ \frac{ q^d \sum_{ \sigma \in \mathcal{L}(\Pi) } \prod_{ j \in \Asc \sigma } q^{ d-j } }{ (1-q) (1-q^2) \cdots (1-q^d) } 
  \ = \ \frac{ q^d \sum_{ \sigma \in \mathcal{L}(\Pi) } q^{ \maj \sigma\op } }{ (1-q) (1-q^2) \cdots (1-q^d) }, 
\end{equation}

\begin{theorem}
Let $G$ be a graph on $d$ vertices. Then
\begin{align*}
  P_G(q)
  \ &= \ \frac{ q^d \sum_{(\rho,\sigma)} q^{ \maj \sigma\op } }{ (1-q) (1-q^2) \cdots (1-q^d) } \\
    &= \ \frac{ q^{ \binom{ d+1 } 2 } \sum_{(\rho,\sigma)} q^{ - \maj \sigma } }{ (1-q) (1-q^2) \cdots (1-q^d) } \, ,
\end{align*}
where each sum ranges over all pairs of acyclic orientations $\rho$ of $G$ and linear extensions $\sigma$ of the poset induced by~$\rho$.
\end{theorem}

\begin{proof}
Since every $G$-partition is a $\Pi_\rho$-partition for exactly one acyclic
orientation $\rho$ of $G$ (and, conversely, every $\Pi_\rho$-partition is a
$G$-partition),
\[
  p_G(n) \ = \sum_{\rho\in \mathcal{A}(G)} p_{ \Pi_\rho } (n)
\]
and so~\eqref{eq:ppartgenfct} gives the first formula:
\[
  P_G(q)
  \ = \ \frac{ q^d \sum_{\rho\in \mathcal{A}(G)} \sum_{ \sigma \in \mathcal{L}(\Pi_\rho) } q^{ \maj \sigma\op } }{ (1-q) (1-q^2) \cdots (1-q^d) } \, .
\]

To see the second formula, we note that each $\rho\in \mathcal{A}(G)$ has a partner
orientation $\overline\rho \in \mathcal{A}(G)$ in which the direction of each edge is
reversed. A linear extension $\sigma \in \mathcal{L}(\Pi_\rho)$ has the corresponding
linear extension $\overline\sigma \in \mathcal{L}(\Pi_{\overline\rho})$ defined via
\[
  \overline\sigma(j)
  \ := \ d+1 - \sigma(d+1-j)
  \  = \ d+1 - \sigma\op(j) \, .
\]
In particular, $j \in \Des(\sigma\op)$ if and only if $j \in \Asc(\overline\sigma)$, and so
\[
  \sum_{(\rho,\sigma)} q^{ \maj \sigma\op }
  \ = \ \sum_{(\overline\rho,\overline\sigma)} q^{ \binom d 2 - \maj \overline\sigma } \, .
\qedhere
\]
\end{proof}

\end{document}
